\documentclass[a4paper, serif, 10pt]{moderncv}

%% moderncv
% style and color
\moderncvstyle{banking}
\moderncvcolor{black}

%% page layout/margins
\usepackage[top=2.5cm, bottom=2.5cm, left=1.5cm, right=1.5cm]{geometry}

\nopagenumbers{}

%% Babel config for Spanish language
% \usepackage[spanish]{babel} has some conflicting issues with moderncv
% so we have to use enable the following options to make it work
%
% ref:
% - https://tex.stackexchange.com/a/140161/36007
\usepackage[spanish,es-lcroman]{babel}

%% fontspec
\usepackage{fontspec}

\IfFontExistsTF{Linux Libertine}{
  \setmainfont[Ligatures={TeX, Common}, Numbers=OldStyle]{Linux Libertine}
}{}
\IfFontExistsTF{Linux Libertine O}{
  \setmainfont[Ligatures={TeX, Common}, Numbers=OldStyle]{Linux Libertine O}
}{}

%% CTeX
% CJK support, used to show CJK characters in the resume
%
% - fontset=none: disable builtin fontset but instead set the CJK font manually
% - heading=false: disable ctex heading
% - punct=kaiming: use kaiming punctuations styles for CJK
% - scheme=plain: use plain scheme, do not override \normalsize font size
% - space=auto: space settings for CJK characters
%
% ref:
% - http://ctan.mirrorcatalogs.com/language/chinese/ctex/ctex.pdf
\usepackage[UTF8, heading=false, punct=kaiming, scheme=plain, space=auto]{ctex}

\IfFontExistsTF{Noto Serif CJK SC}{
  \setCJKmainfont{Noto Serif CJK SC}
}{}
\IfFontExistsTF{Noto Sans CJK SC}{
  \setCJKsansfont{Noto Sans CJK SC}
}{}

%% line spacing
\usepackage{setspace}
\setstretch{1.125}

%% URL styling - use normal text instead of monospace
\urlstyle{same}

% Auto-underline all links
% Defer until \begin{document} so hyperref is already loaded
\AtBeginDocument{%
  \let\oldhref\href
  \renewcommand{\href}[2]{\oldhref{#1}{\underline{#2}}}%
}

%% Patch \httplink and \httpslink to handle full URLs
% The original moderncv macros always prepend http:// or https:// to URLs.
% This causes issues when the URL already has a protocol (e.g., https://example.com)
% resulting in malformed URLs like https://https://example.com.
% This patch checks if the URL already contains "://" and uses it directly if so.
% Note: We use \str_set:Nx (x-expansion) to expand macros like \@homepage first.
\ExplSyntaxOn
\str_new:N \g_yamlresume_url_str

\RenewDocumentCommand{\httpslink}{O{}m}{
  \str_set:Nx \g_yamlresume_url_str {#2}
  \str_if_in:NnTF \g_yamlresume_url_str {://}
    { \url{#2} }
    { \url{https://#2} }
}

\RenewDocumentCommand{\httplink}{O{}m}{
  \str_set:Nx \g_yamlresume_url_str {#2}
  \str_if_in:NnTF \g_yamlresume_url_str {://}
    { \url{#2} }
    { \url{http://#2} }
}
\ExplSyntaxOff

\name{Lucía Fernández Gómez}{}
\title{Gerente de Proyectos Senior | PMP | Transformación Digital}
\phone[mobile]{+34 612 345 678}
\email{lucia.fernandez@email.com}
\homepage{https://www.luciafernandez-pm.com}

\address{Calle de Alcalá 123, 4º B, Madrid, Comunidad de Madrid, España, 28010}{}{}

\extrainfo{{\small \faLinkedin}\ \href{https://www.linkedin.com/in/luciafernandez-pm}{@luciafernandez-pm} {} {} {} • {} {} {} 
{\small \faGithub}\ \href{https://github.com/luciafernandez}{@luciafernandez} {} {} {} • {} {} {} 
{\small \faTwitter}\ \href{https://twitter.com/lucia\_pm}{@lucia\_pm}}

\begin{document}

\maketitle

\section{Información básica}

\cvline{}{Gerente de Proyectos con más de 8 años de experiencia liderando proyectos tecnológicos y de transformación digital en sectores como banca, retail y salud. Especializada en metodologías ágiles (Scrum, Kanban) y en la gestión integral del ciclo de vida del proyecto, desde la definición de requisitos hasta la entrega y el cierre.
%
\begin{itemize}
\item Lideré equipos multidisciplinares de hasta 25 personas, alcanzando una tasa de entrega del 95\% dentro del plazo y presupuesto.
\item Gestioné presupuestos de hasta 2,5 M€, optimizando costes en un 18\% mediante la renegociación con proveedores y la mejora de procesos.
\item Certificada PMP, Scrum Master y ITIL Foundation, con sólida experiencia en gestión de stakeholders y comunicación ejecutiva.
\item Apasionada por la mejora continua y la adopción de herramientas digitales para la gestión de proyectos (Jira, Asana, MS Project).
\end{itemize}}
\section{Educación}

\cventry{sept 2017–jul 2019}
        {Maestría, Dirección de Proyectos, Puntuación: 9.2}
        {Universidad Politécnica de Madrid}
        {\url{https://www.upm.es}}
        {}
        {\begin{itemize}
\item Formación avanzada en dirección y gestión de proyectos complejos
\item Proyecto fin de máster: implantación de PMO en una empresa del sector retail
\end{itemize}
\textbf{Cursos}: Gestión de Proyectos Ágiles, Planificación y Control de Proyectos, Gestión de Riesgos, Liderazgo de Equipos, Finanzas para Gerentes de Proyecto, Metodologías PMBOK}

\cventry{sept 2011–jun 2015}
        {Licenciatura, Administración y Dirección de Empresas, Puntuación: 8.7}
        {Universidad de Sevilla}
        {\url{https://www.us.es}}
        {}
        {\begin{itemize}
\item Base sólida en administración, finanzas y gestión empresarial
\item Participación en programas de intercambio internacional
\end{itemize}
\textbf{Cursos}: Contabilidad Financiera, Dirección Estratégica, Gestión de Operaciones, Comportamiento Organizacional, Marketing}
\section{Experiencia laboral}

\cventry{mar 2021 hasta la fecha}
        {Gerente de Proyectos Senior}
        {Telefónica Tech}
        {\url{https://telefonicatech.com}}
        {}
        {\begin{itemize}
\item Lideré un portfolio de 12 proyectos de transformación digital con un presupuesto total de 2,5 M€.
\item Implementé marcos ágiles híbridos (Scrum + Kanban) que redujeron el time-to-market en un 30\%.
\item Coordiné equipos multidisciplinares de hasta 25 personas distribuidas en 4 países.
\item Gestioné la relación con stakeholders ejecutivos, presentando informes de progreso y métricas KPI.
\item Implanté una PMO que estandarizó procesos y mejoró la trazabilidad de los proyectos.
\end{itemize}
\textbf{Palabras clave}: Scrum, Kanban, Gestión de riesgos, Presupuestos, PMO, Stakeholders}

\cventry{sept 2019–feb 2021}
        {Gerente de Proyectos}
        {Banco Santander}
        {\url{https://www.santander.com}}
        {}
        {\begin{itemize}
\item Dirigí proyectos de desarrollo de software bancario bajo metodología Agile, con equipos de 10-15 personas.
\item Gestioné el ciclo completo de vida de los proyectos: análisis, planificación, ejecución, seguimiento y cierre.
\item Elaboré planes de proyecto, cronogramas y matrices de riesgos, asegurando el cumplimiento de plazos y presupuestos.
\item Facilité ceremonias ágiles (daily stand-ups, sprint planning, retrospectivas) y eliminé impedimentos.
\item Colaboré con product owners y equipos técnicos para priorizar el backlog y entregar valor de forma continua.
\end{itemize}
\textbf{Palabras clave}: Agile, Scrum, Jira, Backlog, Banca, Gestión de proyectos}

\cventry{sept 2015–ago 2017}
        {Analista de Proyectos}
        {Accenture}
        {\url{https://www.accenture.com}}
        {}
        {\begin{itemize}
\item Apoyé la gestión de proyectos de consultoría tecnológica para clientes del sector público y privado.
\item Realicé seguimiento de hitos, elaboración de informes de estado y coordinación de reuniones con clientes.
\item Participé en la definición de alcance, estimación de esfuerzos y control de cambios.
\item Contribuí a la implantación de herramientas de gestión como MS Project y Confluence.
\end{itemize}
\textbf{Palabras clave}: MS Project, Confluence, Consultoría, Informes, Control de cambios}
\section{Idiomas}

\cvline{Español}{Competencia nativa o bilingüe \hfill \textbf{Palabras clave}: Español nativo}
\cvline{Inglés}{Competencia profesional plena \hfill \textbf{Palabras clave}: TOEFL 105, Certificado C1}
\cvline{Francés}{Competencia limitada de trabajo \hfill \textbf{Palabras clave}: DELF B1}
\section{Competencias}

\cvline{Gestión de Proyectos}{Experto \hfill \textbf{Palabras clave}: PMBOK, Planificación, Presupuestos, Gestión de riesgos, Cronogramas}
\cvline{Metodologías Ágiles}{Avanzado \hfill \textbf{Palabras clave}: Scrum, Kanban, Lean, SAFe, Design Thinking}
\cvline{Herramientas}{Avanzado \hfill \textbf{Palabras clave}: Jira, Confluence, MS Project, Asana, Trello, Monday.com}
\cvline{Liderazgo y Comunicación}{Experto \hfill \textbf{Palabras clave}: Liderazgo de equipos, Negociación, Comunicación ejecutiva, Gestión de stakeholders}
\cvline{Análisis y Datos}{Intermedio \hfill \textbf{Palabras clave}: Power BI, Excel avanzado, SQL, Métricas KPI}
\section{Premios}

\cventry{nov 2023}
        {PMI Madrid}
        {Premio a la Excelencia en Gestión de Proyectos}
        {}
        {}
        {Reconocimiento por la dirección del proyecto de transformación digital con mayor impacto en el sector bancario.}

\cventry{dic 2022}
        {Telefónica Tech}
        {Empleada del Año}
        {}
        {}
        {Distinción por liderazgo excepcional y resultados sobresalientes en la gestión de proyectos estratégicos.}
\section{Certificados}

\cventry{jun 2020}
        {Project Management Institute}
        {Project Management Professional (PMP)}
        {\url{https://www.pmi.org/certifications/project-management-pmp}}
        {}
        {}

\cventry{feb 2021}
        {Scrum.org}
        {Professional Scrum Master I (PSM I)}
        {\url{https://www.scrum.org/professional-scrum-master-i-certification}}
        {}
        {}

\cventry{sept 2019}
        {AXELOS}
        {ITIL Foundation}
        {\url{https://www.axelos.com/certifications/itil-certifications}}
        {}
        {}
\section{Publicaciones}

\cventry{may 2022}
        {Guía práctica para la implantación de una PMO ágil}
        {Revista de Dirección de Proyectos}
        {\url{https://www.revistadeproyectos.es/guia-pmo-agil}}
        {}
        {Artículo que describe un marco práctico para implantar una oficina de proyectos ágil en organizaciones medianas.}

\cventry{oct 2021}
        {Métricas de valor en proyectos de transformación digital}
        {Congreso Internacional de Gestión de Proyectos}
        {\url{https://www.congresoproyectos.org/metricas-valor}}
        {}
        {Ponencia sobre indicadores clave para medir el retorno de inversión en iniciativas de transformación digital.}
\section{Proyectos}

\cventry{ene 2022–dic 2022}
        {Desarrollo de una herramienta interna para centralizar la gestión de proyectos y recursos.}
        {Plataforma de Gestión de Proyectos Interna}
        {\url{https://www.telefonicatech.com/proyectos/pm-tool}}
        {}
        {\begin{itemize}
\item Lideré el proyecto desde la definición de requisitos hasta el despliegue en producción.
\item Coordiné un equipo de 8 desarrolladores, 2 diseñadores y 1 analista de negocio.
\item Implementé integraciones con Jira y Power BI para la automatización de informes.
\item Reduje el tiempo de reporting semanal en un 60\%.
\end{itemize}
\textbf{Palabras clave}: Jira, Power BI, Gestión de proyectos, Automatización}

\cventry{mar 2020–ene 2021}
        {Digitalización del proceso de incorporación de clientes para una entidad bancaria.}
        {Transformación Digital del Proceso de Onboarding}
        {\url{https://www.santander.com/onboarding-digital}}
        {}
        {\begin{itemize}
\item Dirigí la implementación de un flujo digital de onboarding con firma electrónica.
\item Reduje el tiempo de alta de clientes de 5 días a 2 horas.
\item Gestioné la coordinación entre equipos de negocio, IT, legal y cumplimiento normativo.
\item Aseguré el cumplimiento de la normativa KYC y GDPR.
\end{itemize}
\textbf{Palabras clave}: Transformación digital, Banca, KYC, GDPR, Onboarding}
\section{Intereses}

\cvline{Deporte}{Senderismo, Natación, Yoga}
\cvline{Voluntariado}{Mentoría, Formación}
\cvline{Lectura}{Novela histórica, Ensayo}
\cvline{Viajes}{Asia, Europa}
\section{Voluntariado}

\cventry{sept 2021–jun 2023}
        {Mentora de Proyectos Sociales}
        {Fundación Integra}
        {\url{https://www.fundacionintegra.org}}
        {}
        {\begin{itemize}
\item Acompaño a organizaciones sin ánimo de lucro en la planificación y ejecución de proyectos sociales.
\item Facilito talleres de gestión de proyectos y metodologías ágiles adaptadas al tercer sector.
\item Ayudo a definir indicadores de impacto y a mejorar la sostenibilidad de las iniciativas.
\end{itemize}}
    

\cventry{ene 2018–dic 2020}
        {Voluntaria en Comité de Eventos}
        {PMI Madrid Chapter}
        {\url{https://pmi-madrid.es}}
        {}
        {\begin{itemize}
\item Colaboré en la organización de congresos y jornadas de gestión de proyectos.
\item Gestioné la logística y coordinación de ponentes para eventos con más de 300 asistentes.
\item Apoyé la comunicación y difusión en redes sociales.
\end{itemize}}
    
\end{document}